% Shared PDF preamble for documents built with pandoc and tectonic (or xelatex).
% bin/make-pdf includes it. Projects with their own build can include it too:
%   --include-in-header=<agent-council>/typeset/pdf-preamble.tex
% and pass the link-colour flags (see "Link colours" below).
%
% Fonts: New Computer Modern (see below). Engine: XeTeX via tectonic.

% ---- Fonts --------------------------------------------------------------
% New Computer Modern, the same family as the deck theme: Sans Bold headings,
% Book body (a sturdier weight of the Computer Modern design), Book math.
% install.sh installs the fonts from typeset/fonts. Under pdflatex, which
% cannot load them, the document falls back to Latin Modern.
\usepackage{iftex}
\ifPDFTeX
  \newcommand\headfont{\sffamily\bfseries}
\else
  \usepackage{fontspec}
  \usepackage{unicode-math}
  % PostScript names: XeTeX resolves these for installed fonts.
  \setmainfont{NewCM10-Book}[ItalicFont=NewCM10-BookItalic, BoldFont=NewCM10-Bold,
    BoldItalicFont=NewCM10-BoldItalic]
  \setmathfont{NewCMMath-Book}
  \newfontfamily\headfont{NewCMSans10-Bold}
\fi
% Heading scale (# 21pt, ## 14pt, ### 12pt, #### body size), with more space
% above than below so each heading attaches to its text. LaTeX's defaults set
% ### at body size, which reads smaller than the body in this sans.
\usepackage{titlesec}
\titleformat{\section}{\headfont\huge}{\thesection}{0.75em}{}
\titleformat{\subsection}{\headfont\Large}{\thesubsection}{0.75em}{}
\titleformat{\subsubsection}{\headfont\large}{\thesubsubsection}{0.75em}{}
\titleformat{\paragraph}{\headfont\normalsize}{\theparagraph}{0.75em}{}
\titlespacing*{\section}{0pt}{3.2ex plus 1ex minus .2ex}{2ex plus .2ex}
\titlespacing*{\subsection}{0pt}{3ex plus 1ex minus .2ex}{1.1ex plus .2ex}
\titlespacing*{\subsubsection}{0pt}{2.6ex plus 1ex minus .2ex}{0.9ex plus .2ex}
\titlespacing*{\paragraph}{0pt}{2.2ex plus 1ex minus .2ex}{0.7ex plus .2ex}

% ---- Characters Latin Modern's text font lacks ---------------------------
% Typed directly in prose (fine in Markdown and HTML), these would print as
% blanks. Draw them from the math font or pifont instead.
\usepackage{iftex}  % \ifPDFTeX: pdflatex (make-pdf default) or XeTeX (fallback)
\usepackage{amssymb}
\usepackage{pifont}
\usepackage{newunicodechar}
\providecommand{\lt}{<}  % MathJax alias for <
\ifPDFTeX
  \newunicodechar{✓}{\ding{51}}
  \newunicodechar{✗}{\ding{55}}
  \newunicodechar{≤}{\ensuremath{\le}}
  \newunicodechar{≥}{\ensuremath{\ge}}
  \newunicodechar{≠}{\ensuremath{\ne}}
  \newunicodechar{≈}{\ensuremath{\approx}}
  \newunicodechar{↔}{\ensuremath{\leftrightarrow}}
  \newunicodechar{⇒}{\ensuremath{\Rightarrow}}
  \newunicodechar{↗}{\ensuremath{\nearrow}}
  \newunicodechar{↘}{\ensuremath{\searrow}}
  \newunicodechar{′}{\ensuremath{^{\prime}}}
  \newunicodechar{″}{\ensuremath{^{\prime\prime}}}
\else
  % XeTeX: unicode-math makes some of these active in math, so draw them by
  % character code from the math font instead of through math mode.
  \newfontfamily\symfont{NewCMMath-Book}
  \newcommand\sym[1]{{\symfont\char"#1\relax}}
  \newunicodechar{✓}{\ding{51}}
  \newunicodechar{✗}{\ding{55}}
  \newunicodechar{≤}{\sym{2264}}
  \newunicodechar{≥}{\sym{2265}}
  \newunicodechar{≠}{\sym{2260}}
  \newunicodechar{≈}{\sym{2248}}
  \newunicodechar{↔}{\sym{2194}}
  \newunicodechar{⇒}{\sym{21D2}}
  \newunicodechar{↗}{\sym{2197}}
  \newunicodechar{↘}{\sym{2198}}
  % Primes (′ ″) are left to unicode-math, which handles them in math. Typed in
  % prose they print blank on this rare fallback path.
\fi
\newunicodechar{⧵}{\ensuremath{\setminus}}
% unicode-math (XeTeX) maps \setminus to U+29F5, which latinmodern-math lacks.
\ifPDFTeX\else\AtBeginDocument{\def\setminus{\mathbin{\backslash}}}\fi

% ---- Tables -------------------------------------------------------------
% A rule under every row, without booktabs' doubled mid and bottom rules.
\usepackage{etoolbox}
\usepackage{longtable,booktabs,array,calc}
\providecommand{\real}[1]{#1}
\makeatletter
\AtBeginEnvironment{longtable}{\footnotesize\let\LT@@cr\LT@tabularcr
  \def\LT@tabularcr{\LT@@cr\hline}\def\midrule{}\def\bottomrule{}}
\makeatother
% Pandoc's columns use \raggedright, which forbids hyphenation. Allow it.
\usepackage{ragged2e}
\let\raggedright\RaggedRight

% ---- Figures ------------------------------------------------------------
\usepackage{graphicx}

% ---- Link colours -------------------------------------------------------
% Pandoc honours coloured links only through its command-line flags:
%   -V colorlinks=true -V linkcolor=xrefblue -V citecolor=citegreen -V urlcolor=xrefblue
% (a \hypersetup here would be overridden by pandoc's hidelinks). This file only
% defines the colours: cross-references and URLs blue, citations green.
\usepackage{xcolor}
\definecolor{xrefblue}{HTML}{1A4D8F}
\definecolor{citegreen}{HTML}{2E7D5B}

% ---- List markers -------------------------------------------------------
% Google Docs progression: filled disc, open circle, filled square, dot.
\usepackage{textcomp}
\renewcommand{\labelitemi}{\textbullet}
\renewcommand{\labelitemii}{\textopenbullet}
\renewcommand{\labelitemiii}{\raisebox{0.15ex}{\rule{0.5ex}{0.5ex}}}
\renewcommand{\labelitemiv}{\textperiodcentered}
